\pdfoutput=1
\documentclass[11pt]{article}
\usepackage[T1]{fontenc}
\usepackage{lmodern}
\usepackage[a4paper,margin=26mm]{geometry}
\usepackage{amsmath,amssymb,amsthm,booktabs,hyperref}
\hypersetup{colorlinks=true,urlcolor=blue,citecolor=blue,
  pdftitle={The ten minimum universal sets of MacMahon cubes},pdfauthor={Jason Davies}}
\newtheorem{theorem}{Theorem}
\newtheorem{lemma}[theorem]{Lemma}
\newtheorem{corollary}[theorem]{Corollary}
\newcommand{\C}{\mathcal C}
\emergencystretch=2em
\newcommand{\lean}[1]{\nolinkurl{#1}}
\makeatletter\g@addto@macro\UrlBreaks{\do\_}\makeatother
\title{The ten minimum universal sets of MacMahon cubes}
\author{Jason Davies\thanks{\href{mailto:jason@jasondavies.com}{\nolinkurl{jason@jasondavies.com}}}}
\date{31 October 2026}
\begin{document}
\maketitle
\begin{abstract}
We classify universal collections of twelve six-coloured MacMahon cubes,
without requiring internal face matching. A row-sum bound in Conway's tableau
forces a symmetric matrix with every row and column sum equal to two.
Corner coverage then forces two disjoint triangles. This proves the
ten-set uniqueness conjecture stated by Johnson and Rold\'an, and also
excludes additional minimum collections containing repeated cube varieties.
The argument is elementary; exhaustive computation is used only as a check
and to obtain additional target-count distributions. The minimum-size theorem,
the complete classification and their discrete geometric interpretation are
formalised in Lean~4.
\end{abstract}

\noindent\textbf{Status and transparency.} This research manuscript and its
Lean formalisation were developed with AI assistance under the author's
direction. The formal proofs have been checked by Lean; independent human
expert review and publication priority beyond the cited sources remain
outstanding. Details of the proof interface and AI involvement appear below.

\section{Conventions and tableau facts}
There are thirty rotational varieties of cubes whose six faces have six
different colours. A collection is \emph{universal} if eight of its physical
cubes can build a $2\times2\times2$ model of every target variety. Interior
face colours need not match. Repeated varieties are allowed in the theorem;
the conjecture in \cite{JR} concerns sets of distinct varieties.

Label Conway's tableau by ordered pairs $ij$ of distinct elements of
$L=\{A,B,C,D,E,F\}$; $ij$ corresponds to the paper's upper/lower-case name
such as $Ab$. Each cube has eight oriented corner types. A corner type is
an ordered triple of face colours modulo cyclic permutation, with a fixed
handedness. There are forty such types.

We use the following finite description of the standard tableau. A corner
type corresponds to a partition of $L$ into two triples, with a directed
3-cycle on each triple. Its six incident cube varieties are precisely the
six directed edges of these cycles. All ten partitions and all four choices
of cycle directions occur. This can be verified from Table~3 of \cite{JR};
Table~\ref{tab:corners} below groups that table by the underlying partition.
The cube $ij$ has exactly the eight corner types containing the arc $i\to j$.
This also gives a compact, purely finite model for the corner incidence.

\begin{table}[ht]
\centering
\begin{tabular}{ll}\toprule
Partition & Four oriented corner types \\\midrule
$ABC\mid DEF$ & 256, 143, 134, 265 \\
$ABD\mid CEF$ & 465, 132, 123, 456 \\
$ABE\mid CDF$ & 345, 126, 162, 354 \\
$ABF\mid CDE$ & 235, 164, 146, 253 \\
$ACD\mid BEF$ & 364, 125, 152, 346 \\
$ACE\mid BDF$ & 356, 142, 124, 365 \\
$ACF\mid BDE$ & 153, 246, 264, 135 \\
$ADE\mid BCF$ & 165, 243, 234, 156 \\
$ADF\mid BCE$ & 136, 254, 245, 163 \\
$AEF\mid BCD$ & 263, 145, 154, 236 \\
\bottomrule\end{tabular}
\caption{The corner types associated with each pair of triples. Directions
are recovered from the six cube names containing the corner in the standard
tableau. Corner numbers use the colour labels of \cite{JR}.}
\label{tab:corners}
\end{table}

Two distinct varieties $ij$ and $uv$ share no corner exactly when
$i=u$, $j=v$, or $(u,v)=(j,i)$. In every other case they share two corners.
Indeed, two distinct arcs in the same directed 3-cycle must form a directed
path; this fixes that cycle, leaving two orientations of the other one.
For two disjoint arcs, the remaining two labels can be assigned to their
respective triples in two ways, each determining both cycle directions.
Common tails, common heads, and reversed arcs cannot coexist in directed
3-cycles. The target itself shares all eight corners.

Consequently each cube is usable for exactly 21 targets. Also, assigning
eight distinct physical cubes to the eight target corner types is both
necessary and sufficient: matching an oriented corner permits a proper
rotation matching the three exposed faces. There are no internal constraints.

\section{Necessity}
Let $a_{ij}$ be the number of copies of variety $ij$ in a collection of size
$K$, with $a_{ii}=0$. Write $r_i=\sum_j a_{ij}$ and $c_j=\sum_i a_{ij}$.

\begin{lemma}[Row bound]
For a universal collection,
\[
 r_i+c_j-2a_{ij}+a_{ji}\le K-8\qquad(i\ne j),
 \qquad 3r_i\le 4K-40,\quad 3c_i\le4K-40.
\]
\end{lemma}
\begin{proof}
For target $ij$, the unusable cubes are those in its row apart from $ij$,
those in its column apart from $ij$, and the copies of $ji$. These classes
are disjoint and total $r_i+c_j-2a_{ij}+a_{ji}$. At least eight cubes must
remain usable, proving the first inequality. Summing over the five
$j\ne i$ gives
\[
5r_i+(K-c_i)-2r_i+c_i=3r_i+K\le5(K-8).
\]
The column bound follows by the analogous sum.
\end{proof}

\begin{lemma}
A universal collection of size twelve consists of one copy of every directed
edge within each of two disjoint triples of tableau labels.
\end{lemma}
\begin{proof}
The row bound gives $r_i\le\lfloor8/3\rfloor=2$. As the six row sums total
twelve, all equal two. The same holds for the columns. If $a_{ij}=0$, the
first inequality now gives $4+a_{ji}\le4$, so $a_{ji}=0$. Thus support is
symmetric.

Consider the undirected support graph on $L$. Every vertex has degree one
or two. If a vertex $i$ has degree one, with neighbour $j$, then $a_{ij}=2$;
column $i$ and support symmetry give $a_{ji}=2$, so this is an isolated
double-copy pair. Every degree-two component is a simple cycle, with one
copy in each direction on every edge.

On six vertices the only possibilities are two triangles, a 6-cycle,
a 4-cycle and a double-copy pair, or three double-copy pairs. All except
two triangles are bipartite with parts of size three. Choose these parts
as the two triples defining a corner type. Every directed edge of that
corner lies within a part, whereas every selected variety joins the parts.
Thus the collection contains no copy of that corner type. A target having
that corner cannot be built, contradicting universality. Only two triangles
remain, and their arc multiplicities are all one.
\end{proof}

\section{Sufficiency and classification}
\begin{lemma}
For any partition $L=P\sqcup Q$ into triples, the twelve directed edges
within $P$ and $Q$ form a universal set.
\end{lemma}
\begin{proof}
Fix target $ij$. Its corners are indexed by a third label $k\notin\{i,j\}$
and either direction of the cycle on the complementary triple. The first
cycle is $i\to j\to k\to i$.

Suppose $i,j\in P$, and write $P=\{i,j,p\}$. For each $k\in Q$, the
complementary cycle contains one arc between the two labels of $Q\setminus
\{k\}$. Assign that cube to the corner. As $k$ and the cycle direction vary,
these are the six distinct directed edges within $Q$. The remaining two
corners have $k=p$; fill them with $jp$ and $pi$, both of which belong to
both corner types. This gives eight distinct cubes.

Now suppose $i\in P$, $j\in Q$, with $P\setminus\{i\}=\{u,v\}$ and
$Q\setminus\{j\}=\{x,y\}$. For the four corners with $k\in\{u,v\}$,
the usable selected varieties are $ki$ and one of $xy,yx$. Their incidence
is that of the four edges of a $K_{2,2}$ with its four vertices. Orient this
4-cycle cyclically to assign a different incident variety to each corner.
For the other four corners, the analogous four varieties are $jx,jy,uv,vu$.
The two groups of four varieties are disjoint, completing the construction.
\end{proof}

\begin{theorem}
The minimum size of a universal collection is twelve. Exactly ten labelled
collections attain that size, even when repeated varieties are permitted.
They are the two-triangle sets described above. In particular, Conjectures
4 and 8 of Johnson and Rold\'an hold.
\end{theorem}
\begin{proof}
Each physical cube is usable for 21 targets, so universality requires
$21K\ge30\cdot8=240$, hence $K\ge12$. This lower bound is due to
Haraguchi \cite{H}. The preceding lemmas give necessity and sufficiency at
twelve. There are $\frac12\binom63=10$ unordered partitions into triples.
\end{proof}

\section{Computational checks and additional results}
An exhaustive enumeration of all $\binom{30}{12}=86,493,225$ subsets returns
exactly the ten sets above. The implementation also reproduces the published
distribution for all $\binom{30}{8}=5,852,925$ eight-cube sets.
Two different feasibility algorithms are used: bipartite matching, and a
graph-component test. In the latter, every usable non-target cube is an
edge joining its two compatible target corners. A connected tree component
has deficit one; a component containing a cycle has deficit zero. Thus a
target is buildable precisely when the number of tree components does not
exceed the number of available target copies (zero or one in the subset
enumerations). The table-based implementation precomputes all $2^{20}$ edge
subsets for each target. No symmetry reduction is used in the exhaustive
subset loops.

Independent geometric code generates all 24 proper cube rotations, all
30 varieties and all 40 corner types from face-colour permutations. It checks
the tableau, the directed-triangle correspondence, and explicit physical
placements for all 300 set--target pairs. Independently enumerating the 70
simple 2-regular graphs gives 60 six-cycles, each building 18 targets, and
the ten universal triangle-pair sets. The ten sets form one colour-permutation
orbit, with stabilizer order 72.

For sets of distinct varieties the exhaustive target-count ranges are:
\[
\begin{array}{c|rrrrrr}
\text{set size}&8&9&10&11&12&13\\\hline
\text{minimum targets}&0&0&1&2&8&10\\
\text{maximum targets}&5&9&20&22&30&30
\end{array}
\]
No set of size at most twelve builds exactly 29 targets. At size thirteen,
827,250 sets build exactly 29 targets and 53,100 are universal.
These enumeration results address further questions in Section 6 of
\cite{JR}; unlike the uniqueness theorem, their upper bounds rely on the
complete computational enumeration. Full histograms, example masks,
source code and independent checkers accompany this draft.

\section{Formal verification and reproducibility}
The Lean~4 development accompanies this manuscript at
\url{https://github.com/jasondavies/macmahon-min-universal-sets}.
It is pinned to Lean~4.35.0 and mathlib~v4.35.0. The main results are
\lean{MacMahon.minimum_size}, \lean{MacMahon.physical_exactly_ten},
and \lean{MacMahon.physical_iff_two_triangles}.

An inventory is a natural-number multiplicity for each ordered pair of
six tableau labels, with zero diagonal. A construction is an injective
assignment of eight physical copies to the eight positions, each with a
proper rotation matching its three exposed faces. The formalisation proves
that the rotation table is exactly the set of determinant-positive
signed-axis permutations, and that the thirty representatives and their
rotations give all 720 assignments of six distinct colours to six faces,
without duplication. It also proves the equivalence of corner matching
and exposed-face matching.

The inequalities and symmetry argument are proved for arbitrary natural
multiplicities. The remaining classification is checked through 130
arithmetic proof branches, split across fifteen modules. All finite checks
use kernel reduction, rather than a trusted external computation. Explicit
witnesses certify all 300 minimum-set--target constructions. The proof
library contains no added axioms or unfinished proofs; guarded audits allow
only Lean's standard propositional extensionality, classical choice and
quotient axioms. The larger target-count enumerations above are not included
in these formal theorems.

For Palomar verification, \texttt{Challenge.lean} imports only mathlib and
states ten advertised results with intentional proof holes, independently
of the proof library. \texttt{Solution.lean} imports the completed proofs.
The comparator configuration checks equality of the statements, their
permitted axioms, and kernel acceptance. Local verification outputs and
reproduction commands are recorded in the repository; a successful local
check is distinct from registry acceptance.

\paragraph{AI disclosure.}
OpenAI Codex assisted with the investigation, compiled solvers, geometric
checks, proof argument, Lean development, manuscript and submission packaging,
under the author's direction. No claim of independent human expert review
is made here. The proof generators are not trusted components: their output
must be accepted by Lean's kernel.

\begin{thebibliography}{9}
\bibitem{JR} I.~Johnson and E.~Rold\'an,
\emph{Solution Numbers for Eight Blocks to Madness Puzzle},
arXiv:2407.13208v1 (2024). \url{https://arxiv.org/abs/2407.13208}.
\bibitem{H} K.~Haraguchi,
\emph{On a generalization of ``Eight Blocks to Madness'' puzzle},
Discrete Mathematics \textbf{339} (2016), 1400--1409.
Preprint arXiv:1408.3696v2. \url{https://arxiv.org/abs/1408.3696}.
\end{thebibliography}
\end{document}
